% Options for packages loaded elsewhere
\PassOptionsToPackage{unicode}{hyperref}
\PassOptionsToPackage{hyphens}{url}
\PassOptionsToPackage{dvipsnames,svgnames,x11names}{xcolor}
\documentclass[
  10pt,
  fontset=fandol]{ctexart}
\usepackage{xcolor}
\usepackage[margin=16mm]{geometry}
\usepackage{amsmath,amssymb}
\setcounter{secnumdepth}{-\maxdimen} % remove section numbering
\usepackage{iftex}
\ifPDFTeX
  \usepackage[T1]{fontenc}
  \usepackage[utf8]{inputenc}
  \usepackage{textcomp} % provide euro and other symbols
\else % if luatex or xetex
  \usepackage{unicode-math} % this also loads fontspec
  \defaultfontfeatures{Scale=MatchLowercase}
  \defaultfontfeatures[\rmfamily]{Ligatures=TeX,Scale=1}
\fi
\usepackage{lmodern}
\ifPDFTeX\else
  % xetex/luatex font selection
\fi
% Use upquote if available, for straight quotes in verbatim environments
\IfFileExists{upquote.sty}{\usepackage{upquote}}{}
\IfFileExists{microtype.sty}{% use microtype if available
  \usepackage[]{microtype}
  \UseMicrotypeSet[protrusion]{basicmath} % disable protrusion for tt fonts
}{}
\makeatletter
\@ifundefined{KOMAClassName}{% if non-KOMA class
  \IfFileExists{parskip.sty}{%
    \usepackage{parskip}
  }{% else
    \setlength{\parindent}{0pt}
    \setlength{\parskip}{6pt plus 2pt minus 1pt}}
}{% if KOMA class
  \KOMAoptions{parskip=half}}
\makeatother
\setlength{\emergencystretch}{3em} % prevent overfull lines
\providecommand{\tightlist}{%
  \setlength{\itemsep}{0pt}\setlength{\parskip}{0pt}}
\usepackage{amsmath,amssymb,mathtools,needspace}
\xeCJKDeclareCharClass{CJK}{"2032,"0400->"04FF,"2160->"216B,"2460->"2473,"25A0,"25C7,"25CB,"25CF}
\IfFontExistsTF{Arial Unicode MS}{\xeCJKsetup{AutoFallBack=true}\setCJKfallbackfamilyfont{\CJKrmdefault}{Arial Unicode MS}}{}
\setcounter{secnumdepth}{0}
\setlength{\emergencystretch}{2em}
\allowdisplaybreaks
\usepackage{bookmark}
\IfFileExists{xurl.sty}{\usepackage{xurl}}{} % add URL line breaks if available
\urlstyle{same}
\hypersetup{
  pdftitle={多项式求值的秦九韶算法},
  colorlinks=true,
  linkcolor={Maroon},
  filecolor={Maroon},
  citecolor={Blue},
  urlcolor={Blue},
  pdfcreator={LaTeX via pandoc}}

\title{多项式求值的秦九韶算法}
\author{}
\date{}

\begin{document}
\maketitle

\subsubsection{6.5.1 多项式求值的秦九韶算法}\label{ux591aux9879ux5f0fux6c42ux503cux7684ux79e6ux4e5dux97f6ux7b97ux6cd5}

多项式的一个重要特点是求值方便。

设给定多项式

\[
f(x)=a_0x^n+a_1x^{n-1}+\cdots+a_{n-1}x+a_n,
\]

其中系数 \(a_i\ (0\leqslant i\leqslant n)\) 均为实数。用一次式 \(x-x_0\) 除 \(f(x)\)，商记作 \(P(x)\)，余数显然等于 \(f(x_0)\)，即有

\[
f(x)=f(x_0)+(x-x_0)P(x).
\tag{6.5.1}
\]

现在具体确定 \(P(x)\) 与 \(f(x_0)\)。令

\[
p(x)=b_0x^{n-1}+b_1x^{n-2}+\cdots+b_{n-2}x+b_{n-1},
\]

代入式（6.5.1），并比较两端同次幂的系数，得

\[
\begin{cases}
a_0=b_0,\\
a_i=b_i-x_0b_{i-1},\quad1\leqslant i\leqslant n-1,\\
a_n=f(x_0)-x_0b_{n-1}.
\end{cases}
\]

从而有

\[
\begin{cases}
b_0=a_0,\\
b_i=a_i+x_0b_{i-1},\quad1\leqslant i\leqslant n,\\
f(x_0)=b_n.
\end{cases}
\tag{6.5.2}
\]

\href{../../source-images/pdf-172.jpeg}{原页}

这里提供的一种计算函数值 \(f(x_0)\) 的有效算法称为秦九韶法①。这种算法的优点是计算量小，结构紧凑，容易编制计算程序。

进一步考察 \(f(x)\) 的 Taylor 展开式

\[
f(x)=f(x_0)+f'(x_0)(x-x_0)+\frac{f''(x_0)}{2!}(x-x_0)^2+\cdots+\frac{f^{(n)}(x_0)}{n!}(x-x_0)^n.
\]

将它表示为式（6.5.1）的形式，则式中

\[
P(x)=f'(x_0)+\frac{f''(x_0)}{2!}(x-x_0)+\cdots+\frac{f^{(n)}(x_0)}{n!}(x-x_0)^{n-1}.
\]

由此可见，导数 \(f'(x_0)\) 又可看作 \(P(x)\) 用因式 \(x-x_0\) 相除得出的余数，即有

\[
P(x)=f'(x_0)+(x-x_0)Q(x),
\]

其中 \(Q(x)\) 是 \(n-2\) 次多项式。设

\[
Q(x)=c_0x^{n-2}+c_1x^{n-3}+\cdots+c_{n-3}x+c_{n-2},
\]

那么，再用秦九韶算法又可求出值 \(f'(x_0)\)，对应于式（6.5.2），这里计算公式是

\[
\begin{cases}
c_0=b_0,\\
c_i=b_i+x_0c_{i-1},\quad1\leqslant i\leqslant n-1,\\
f'(x_0)=c_{n-1}.
\end{cases}
\tag{6.5.3}
\]

继续这一过程，可以依次求出 \(f(x)\) 在点 \(x_0\) 的各阶导数。

\end{document}
